% !TeX root = ../navier-stokes-manuscript.tex
\Needspace{12\baselineskip}
\subsection{What the missing rectangle estimate would give}\label{sub:BodyCompatibleIntersection}

Assume the missing estimate, without claiming its proof.
\begin{assumption}[Whole-terminal rectangle producer]\label[assumption]{thm:BodyWholeTerminalRectangleTheorem}
On the finite-lifespan branch, every fixed finite rectangle satisfies
\begin{equation}\label{eq:BodyWholeTerminalRectangleEstimate}
  \sup_{0<T<T_*}\mathfrak R_{N,M}(u;T)<\infty,
  \qquad N,M\in\Nzero.
\end{equation}
The bound may depend on \(N\) and \(M\); it is uniform only in the cutoff \(T<T_*\).
\end{assumption}

Fix \(N\) and \(M\).
Every term in \eqref{eq:BodyRectangleNorm} is the integral over \((0,T)\) of a nonnegative squared norm.
As \(T\uparrow T_*\), the intervals \((0,T)\) exhaust \(I_*=(0,T_*)\), and
\[
  \int_{I_*}f(t)\,dt
  =
  \lim_{T\uparrow T_*}\int_0^T f(t)\,dt
\]
for each of those integrands.
Under \Cref{thm:BodyWholeTerminalRectangleTheorem}, the limit is finite for every row in the fixed rectangle:
\begin{equation}\label{eq:BodyWholeTerminalAdjacentRows}
  V_n\in L^2(I_*;H^{M+1}),
  \qquad
  V_{n+1}=\partial_tV_n\in L^2(I_*;H^{M-1}),
  \qquad 0\leq n\leq N.
\end{equation}
The adjacent trace identity now carries every selected row continuously to the
terminal time:
\[
  V_n\in C([0,T_*];H^M),
  \qquad 0\leq n\leq N.
\]

Under the same assumption, \Cref{thm:WholeTerminalCompatibleRectangleEstimate} applies the pressure recurrence, Riesz-transform bounds, and finite-order Sobolev products to give
\begin{equation}\label{eq:BodyWholeTerminalPressureRows}
  Q_n\in L^1(I_*;H^M),
  \qquad 0\leq n\leq N.
\end{equation}
The velocity and pressure coordinates in
\eqref{eq:BodyWholeTerminalAdjacentRows}--\eqref{eq:BodyWholeTerminalPressureRows}
form the whole-interval rectangle \(\cR_{N,M}[u_0]\).

Place that rectangle in
\[
  \mathscr X_{N,M}
  :=
  \prod_{n=0}^{N}L^2(I_*;H^{M+1})
  \times
  \prod_{n=0}^{N}L^2(I_*;H^{M-1})
  \times
  \prod_{n=0}^{N}L^1(I_*;H^M),
\]
whose three factors carry \((V_n)_{n=0}^{N}\),
\((V_{n+1})_{n=0}^{N}\), and \((Q_n)_{n=0}^{N}\).

Now choose \(N'\leq N\) and \(M'\leq M\).
The restriction
\(\rho_{N',M'}^{N,M}:\mathscr X_{N,M}\to\mathscr X_{N',M'}\)
discards the temporal rows above \(N'\) and uses
\[
  H^{M+1}\hookrightarrow H^{M'+1},
  \qquad
  H^{M-1}\hookrightarrow H^{M'-1},
  \qquad
  H^M\hookrightarrow H^{M'}
\]
on the retained coordinates.
Nothing is reconstructed on the overlap.
Each retained coordinate is the same distributional derivative of the same velocity or normalized pressure, viewed at the lower Sobolev height.
Hence
\begin{equation}\label{eq:BodyRectangleRestrictionCompatibility}
  \rho_{N',M'}^{N,M}\cR_{N,M}[u_0]
  =\cR_{N',M'}[u_0].
\end{equation}

Under the rectangle assumption, the whole-interval family is therefore
\begin{equation}\label{eq:DatumSpecificCompatibleIntersection}
  \cI[u_0]
  :=\bigl(\cR_{N,M}[u_0]\bigr)_{N,M\in\Nzero}
  \in\varprojlim_{N,M}\mathscr X_{N,M}.
\end{equation}
The limit records exact agreement under restriction.
It is not convergence in one norm over all derivative orders, and it does not create a second velocity field.

\begin{proposition}[Conditional mixed intersection]\label[proposition]{prop:BodyCompatibleProjectiveAssembly}
Under \Cref{thm:BodyWholeTerminalRectangleTheorem}, the velocity component of
\(\cI[u_0]\) satisfies
\begin{equation}\label{eq:BodyProducedMixedIntersection}
  u\in
  \bigcap_{r=0}^{\infty}
  \bigcap_{m=0}^{\infty}
  H^r(I_*;H^m(\R^3)).
\end{equation}
In particular,
\begin{equation}\label{eq:CompleteSpatialRow}
  u\in\bigcap_{m=0}^{\infty}L^2(I_*;H^m(\R^3)).
\end{equation}
\end{proposition}

\begin{proof}
Fix finite \(r\) and \(m\), then choose \(N\geq r\) and \(M\geq m\).
The upper rows of that rectangle give
\[
  \partial_t^nu\in L^2(I_*;H^{M+1})
  \hookrightarrow L^2(I_*;H^m),
  \qquad 0\leq n\leq r.
\]
These are the distributional time derivatives of the same velocity field, so
\(u\in H^r(I_*;H^m)\).
Since \(r\) and \(m\) were arbitrary, all finite coordinates belong to the intersection.
Taking \(r=0\) gives \eqref{eq:CompleteSpatialRow}.
\end{proof}

The cutoff-uniform rectangle bounds are sufficient for the mixed intersection.
No argument in this subsection derives those bounds from \((u_0,\nu)\).
The datum-to-rectangle step remains open, so the intersection and every endpoint or continuation consequence that uses it remain conditional.
